\clearpage
\VerdierPhysicalPageAnchor{0108}%
\begingroup
\small
\setlength{\abovedisplayskip}{4pt plus 2pt minus 2pt}%
\setlength{\belowdisplayskip}{4pt plus 2pt minus 2pt}%

\VerdierSemanticAnchor{ch1-app-signs}
\section*{Appendix: Comments on the Choice of Signs}
\addcontentsline{toc}{section}{Appendix: Comments on the Choice of Signs}

\VerdierSemanticAnchor{ch1-app-signs-cone}
The cone functor is determined up to unique isomorphism by its adjunction
properties (3.1.3). In order to associate easily to a morphism of
$\operatorname{tr}_1(C)$ a triangle of $\operatorname{comp}(C)$, we chose as
translation functor the functor “cone of the morphism
$X\longrightarrow 0$” ((3.1.7) and (3.2.2.7)), i.e. shifting the degrees and
changing the sign of the differential.

\VerdierSemanticAnchor{ch1-app-signs-strict}
Let $F\colon\prod_{i\in[n]}C_i\longrightarrow C$ be a multiadditive functor.
The category $\prod_{i\in[n]}\operatorname{comp}(C_i)$ is naturally endowed
with the structure of a $\mathbb Z[n]$-category. If one wants the functor
\[
  F^{[n]}\colon\prod_{i\in[n]}\operatorname{comp}(C_i)
  \longrightarrow\operatorname{comp}^{n}(C)
\]
to be a strict $\mathbb Z[n]$-functor, one is led to the choice of
conventions that we made for translations: shifting the degrees in one
direction and changing the sign of the corresponding differential.

\VerdierSemanticAnchor{ch1-app-signs-simple-complex}
We adopted the conventions of [1] for the choice of the associated
simple-complex functor (except that, in our multiple complexes, the
differentials commute with one another). This makes it possible to define the
functor $K_{\Pi}F\colon\prod_i K(C_i)\longrightarrow K(C)$. It then turns out
that $K_{\Pi}F(X_1,\ldots,X_i[1],\ldots,X_n)$ is different from
$K_{\Pi}F(X_1,\ldots,X_i,\ldots,X_n)[1]$. In order nevertheless to assert
that the functor $K_{\Pi}F$ sends distinguished triangles to distinguished
triangles, one is led to define the isomorphisms

\VerdierSemanticAnchor{ch1-app-signs-shift-isomorphism}
\[
  p_i\colon K_{\Pi}F(X_1,\ldots,X_i[1],\ldots,X_n)
  \xrightarrow{\ \sim\ }
  K_{\Pi}F(X_1,\ldots,X_i,\ldots,X_n)[1]
\]
(proof of (3.4.3)).

\VerdierSemanticAnchor{ch1-app-signs-anticommutativity}
It then turns out that, for $i\ne j$, one has anticommutative diagrams:

\VerdierSemanticAnchor{ch1-app-signs-diagram}
\begin{center}
  \resizebox{\textwidth}{!}{%
    $\begin{tikzcd}[ampersand replacement=\&,column sep=3.2em,row sep=4.3em]
      K_{\Pi}F(X_1,\ldots,X_i[1],\ldots,X_j[1],\ldots)
        \arrow[r,"p_j"]
        \arrow[d,"p_i"']
        \& K_{\Pi}F(X_1,\ldots,X_i[1],\ldots,X_j,\ldots)[1]
          \arrow[d,"p_i"] \\
      K_{\Pi}F(X_1,\ldots,X_i,\ldots,X_j[1],\ldots)[1]
        \arrow[r,"p_j"]
        \& K_{\Pi}F(X_1,\ldots,X_i,\ldots,X_j,\ldots)[2].
    \end{tikzcd}$%
  }
\end{center}

\VerdierSemanticAnchor{ch1-app-signs-koszul-choice}
We then showed in (1.5.2) how to extend these isomorphisms coherently to
isomorphisms $p(\sigma)$ defined for every $\sigma\in\mathbb Z[n]$. An
arbitrary choice enters here. It consists in choosing a cocycle in a
cohomology class that is itself well determined: the Koszul class (1.6.3).
In order to conform to commonly accepted rules (for example [1]), we chose in
every case the Koszul cocycle (1.6.6).
\endgroup
